% Source: https://artofproblemsolving.com/community/c3h1962913_small_red_squares
% Stable identity: aops-topic-1962913-q01
No. Consider the total perimeter $L$ of the red region, that is, the number of unit edges that separate a red square from a non-red square or from the outside of the board.

When a white square with at least two red neighbours turns red, at least two of its four sides were edges of the red boundary; they stop being boundary edges, and at most the other two sides become new boundary edges. Hence $L$ never increases. (If several squares turn red in the same step, colour them one at a time: each still has at least two red neighbours at its turn, so each single change does not increase $L$.)

At the start there are $10$ red squares, so $L\le 4\cdot10=40$. If the whole board were red, its perimeter would be $4\cdot16=64>40$. Hence the whole board can $\boxed{\text{never}}$ become red.
